\scanpage{75}
\begin{center}\fbox{Fonctions elliptiques}\end{center}

\inserted{\uline{Expression \unsure{des|en} fonctions elliptiques par $\unsure{\wp|p}u,\wp' u$}}

1) \uuline{Généralités}. On appelle \uline{fonction elliptique} une fonction méromorphe $f$ dans $C$ qui admet deux périodes $\omega_1,\omega_2$ linéairement ind. sur $R$ ; si $f$ n'est pas constante, son groupe des périodes, (sous-groupe fermé de $C=R^2$) est discret, soit $\unsure{\mathfrak P|\mathfrak B}$\inserted{ et de rang 2.} On étudie \inserted{l'ensemble} des fonctions elliptiques qui admettent $\mathfrak P$ comme groupe des périodes.

\deleted{\unsure{Soit donc|Soit la}} Soit $\{2\omega_1,2\omega_2\}$ une base de $\mathfrak P$ sur $Z$ (définie à une substitution entière de déterminant $\pm1$ près) on étudie l'ens des fonctions méromorphes telles que
\[
f(z+2\omega_1)=f(z+2\omega_2)=f(z)\tag*{(1)}
\]
Il forme un \uline{corps} $E$, qu'on peut évidemment identifier au corps des fonctions méromorphes sur la variété holomorphe quotient $C/\mathfrak P$, -- qui est topologiquement isomorphe au tore. -- On sait, d'après \deleted{\unsure{le|la}} \unsure{théorie|théorème} \deleted{\unsure{finitiste|finitude|générale}} \inserted{de Riemann}, que \deleted{\unsure{pour}} \inserted{pour} \deleted{toute} $f\in E$ non constante \deleted{\unsure{on}} (et il en existe) \deleted{\illegible{1}} $E$ est de degré fini sur $C(f)$, i. e \deleted{\unsure{toute}} fonction elliptique \inserted{$g\in E$} est \deleted{\unsure{égale}} solution d'une équation $P(f,g)=0$, $P$ polynôme en $f$ et $g$ \inserted{etc.} On peut ici obtenir élémentairement des résultats plus précis.

a') \deleted{\unsure{Existence d'une fon}} \uline{Une fonction elliptique sans pôle est une constante}

-- \uline{immédiat par} Liouville

\margin{\unsure{a)|a')} Somme des résidus nulle ; nb de zéros \unsure{égale|égal au} nb \unsure{des|de} pôles. \unsure{corollaire} : au moins deux \unsure{distincts|d'ordre}.}

b) Existence d'une fonction elliptique \unsure{non constante}
\[
\wp(z)=\frac{1}{z^2}+\sum_{m,n}'\left(\frac{1}{(z-\unsure{m|2m}\omega_1\unsure{+|-}n\omega_2)^2}-\frac{1}{(m\omega_1+n\omega_2)^2}\right)\tag*{(2)}
\]
définit une fonction méromorphe dans \deleted{\unsure{tout}} le plan, ayant \deleted{pour pôles} $\mathfrak P$ comme ensemble des pôles (pôles doubles)

On a
\[
\wp'(z)=-\frac{2}{z^3}-2\sum_{m,n}'\frac{1}{(z\unsure{+|-}m\omega_1+n\omega_2)^3}
=-2\sum\frac{1}{(z-2m\omega_1\unsure{-|+}2n\omega_2)^3}\tag*{3)}
\]

% archive sheet number: 74
